\documentclass[10pt,a4paper]{amsart}
\usepackage[margin=19mm]{geometry}
\usepackage{amsmath,amssymb,mathtools}
\usepackage[scr=rsfs,cal=euler]{mathalfa}
\usepackage{xfrac}
\usepackage[unicode,hidelinks,bookmarks=false]{hyperref}
\newcommand{\ie}{i.e.\ }
\newcommand{\cf}{cf.\ }
\newcommand{\resp}{respectively\ }
\newcommand{\Subsection}{\subsection}
\providecommand{\nobreakdash}{\nobreak\hbox{-}}
\setlength{\emergencystretch}{2em}
\multlinegap0pt
\usepackage{bm}
\usepackage{bbm}
\usepackage{dsfont}
\usepackage{xspace}
\usepackage{cancel}
\usepackage{mathtools}
\usepackage{amsthm}
\makeatletter
\@ifundefined{Theorem}{\newtheorem{Theorem}{Theorem}}{}
\@ifundefined{Corollary}{\newtheorem{Corollary}[Theorem]{Corollary}}{}
\@ifundefined{Definition}{\newtheorem{Definition}[Theorem]{Definition}}{}
\makeatother
\newcommand{\R}{\mathbb{R}}
\newcommand{\abs}[1]{\left|#1\right|}
\newcommand{\del}{\delta}
\newcommand{\de}{\mathrm{d}}
\newcommand{\e}{\varepsilon}
\newcommand{\mr}{\mathrm}
\newcommand{\sgn}{\operatorname{sgn}}
\title[Homogenization in one-dimensional higher-order non-local mod]{Homogenization in one-dimensional higher-order non-local models of phase transitions}
\author{Fabrizio Caragiulo, Sergio Scalabrino, Edoardo Voglino}
\date{Source: arXiv:2604.12689v1 (CC BY 4.0)}
\begin{document}
\maketitle

\noindent\emph{Source and licence.} Homogenization in one-dimensional higher-order non-local models of phase transitions. Version arXiv:2604.12689v1 (CC BY 4.0), distributed under the Creative Commons Attribution 4.0 International licence, as recorded in the arXiv licence metadata for this exact version and shown on the abstract page https://arxiv.org/abs/2604.12689v1. Full rights notice: licenses/arxiv-2604.12689-CC-BY-4.0.txt.\par\smallskip

\noindent\textit{Editorial scope.} Counted unit: the Theorem whose source label is \texttt{ineq:limsup} in the article (source0.tex lines 1190--1197, source label \texttt{ineq:limsup}), delivered together with the article's own complete original proof of it (source0.tex lines 1198--1319, one \texttt{proof} environment of 122 lines). No mathematics is written, repaired or re-derived: statement and proof are the article's own text line for line. Required context retained uncounted: def:transition\_energies (lines 314--333); coroll:min\_equivalence (lines 885--888). Every definition, display and label of those lines is retained unchanged. Omitted from this package and not redistributed: 1--313, 334--884, 889--1189, 1320--1634. Nothing inside a retained excerpt is omitted or altered: every retained body is delivered exactly as the article's lines are written, so no omission receipt of any kind accompanies this package. Citations: the retained excerpts carry no citation command at all, so no bibliography entry of the article is transcribed and no reference is redistributed. Reference adaptation: none was needed and none was made; every reference inside the delivered unit resolves inside it, so no unresolved reference or citation is produced. Printed slips kept literal: no printed word, symbol, hypothesis or number of the counted statement or of its proof was altered. Layout and class: the article's own class is \texttt{article} with packages including mathtools, amssymb, amsthm, esint, empheq, etoolbox, mathrsfs, graphicx, tikz, cancel, hyperref, biblatex; the wrapper is the lane's pinned \texttt{amsart} editorial wrapper at 10pt on A4, which restates the theorem-like environments the retained text needs and adds the article's own macro definitions that the retained text needs (\texttt{\textbackslash R}, \texttt{\textbackslash abs}, \texttt{\textbackslash de}, \texttt{\textbackslash del}, \texttt{\textbackslash e}, \texttt{\textbackslash mr}, \texttt{\textbackslash sgn}), with extra wrapper packages (bm, bbm, dsfont, xspace, cancel, mathtools). The delivered numbering is the unit's own. Assets: no retained span contains a figure environment, a graphics-inclusion command, a PDF-page inclusion, a table or a TikZ picture; no image file is copied into this package, the package declares no asset dependency and no image file is redistributed with it. Licence: version arXiv:2604.12689v1 of this article is distributed under the Creative Commons Attribution 4.0 International licence, as recorded in the arXiv licence metadata for this exact version and shown on the abstract page https://arxiv.org/abs/2604.12689v1. A full rights notice is delivered at licenses/arxiv-2604.12689-CC-BY-4.0.txt. Characters: every retained excerpt is pure ASCII, so the delivered engine, XeLaTeX with its default Latin Modern text font, reports no missing character.
\begin{Definition}[Transition energies]
\label{def:transition_energies}
Given a homogenization kernel $a$ and $\omega\in \{\pm 1\}$, we define the \emph{transition energies} at finite and infinite length respectively as
    \begin{align*}
        m^\omega(a,T) &\coloneqq \inf \big\{\Phi^{a}(v)\;\big|\; v\in H^{k+s}_{\operatorname{loc}}(\R), \; v(x) = \omega \sgn(x)  \text{ for } \abs{x} > T\big\},\\
         m^{\omega}(a)  &\coloneqq \inf \big\{\Phi^{a}(v)\;\big|\; v\in H^{k+s}_{\operatorname{loc}}(\R), \; \lim_{x\to \pm\infty}v(x) = \pm \omega \big\}.
    \end{align*}
We also introduce the notations
\begin{equation*}
    m_{\lambda,r}^{\omega}(a)\coloneqq m^\omega\Big(a\Big(\frac{\cdot+r}{\lambda},\frac{\cdot+r}{\lambda}\Big)\Big), \quad m_\lambda^{\omega}(a)\coloneqq m^\omega\big(a(\cdot\,/\lambda,\, \cdot\,/\lambda)\big),
\end{equation*}
 for $\lambda \in (0,+\infty)$, and
 \begin{equation*}
\begin{split}
    m^\omega_0(a)=m^\omega(\bar a), \quad  &\bar a=\int_{(0,1)^2}a(x,y)\,\de x\,\de y,\\
    m^\omega_\infty(a)=m^\omega(a_{\inf}),\quad  &a_{\inf}=\inf_{t\in \R}\{a(t,t)\}.
\end{split}
\end{equation*}
The $T$-dependent versions $m_{\lambda,r}^{\omega}(a, T)$ and so on are defined analogously, by considering $m^\omega(\cdot\,, T)$ at the right-hand sides.
\end{Definition}

\begin{Corollary}\label{coroll:min_equivalence}
For any homogenization kernel $a$ and $\omega\in \{\pm 1\}$
    \[m^{\omega}(a) = \lim_{T\rightarrow +\infty} m^{\omega}(a,T) = \inf_{T>0} m^{\omega}(a,T).\]
\end{Corollary}

\begin{Theorem}
    For every $u \in BV((0,1), \{\pm 1\})$ there exists a sequence $(u_\e)_{\e}$ in $H^{k+s}((0,1))$ such that $u_\e \to u$ in measure as $\e \to 0$ and, for any $\lambda\in [0,+\infty)$,
    \begin{equation}
    \label{ineq:limsup}
        \limsup_{\substack{\e\to 0 \\ \del/\e\to\lambda}} F_{\e,\del}(u_\e) \leq m^+_\lambda\#S^+(u) + m^-_\lambda\# S^-(u).
    \end{equation}
    If we further suppose $a$ to be continuous, then \eqref{ineq:limsup} is true also for $\lambda=\infty$.
\end{Theorem}

\begin{proof}
    Let us enumerate the elements of the jump-set,
    \[S(u)= \{t_1,\dots, t_M\},  \quad 0\eqqcolon t_0< t_1<t_2<\ldots<t_M<t_{M+1}\coloneqq 1.\] 
    and let  $2\displaystyle\tau=\min_{i\ne j}|t_i-t_j|$ be the minimal distance among jump points or between a jump point and an edge of the interval. We also let
    \[s_j = \sgn\big(u(t_j^+)- u(t_j^-)\big),\quad j=1,\ldots, M\] 
        be the ``direction" of the jumps. L
        
    By Corollary \ref{coroll:min_equivalence}, for any $\lambda\in [0,+\infty]$, \[m^\pm_{\lambda} = \inf_T m^\pm_{\lambda}(T).\]    
   Consequently, for any $N>0$ and $\lambda\in [0,+\infty]$, there exist $T>0$ and a function $v^\pm_\lambda\in H^{k+s}_{\operatorname{loc}}(\R)$ such that $v^\pm_\lambda(x) \equiv \pm\sgn(x)$ for $\abs{x}>T$ and, recall Definition \ref{def:transition_energies},
    \begin{equation}\label{eq:approx_choice}
       \begin{split}
            \Phi^{a(\cdot\,/\lambda)}(v^\pm_\lambda) &\leq m_\lambda^{\pm} + \frac{1}{N}, \quad \text{for }\lambda\in (0,+\infty),\\
             \Phi^{\bar a}(v^\pm_0) &\leq m_0^{\pm} + \frac{1}{N},\\
              \Phi^{a_{\inf}}(v^\pm_\infty) &\leq m_\infty^{\pm} + \frac{1}{N}.\\
       \end{split}
    \end{equation} 
    Moreover, it is also possible to suppose that $T>1$ and, when $\lambda \ne +\infty$, even that $T>\lambda$.
    We note that by construction, $||v^\pm_\lambda||_{L^\infty(\R)} =  ||v^\pm_\lambda||_{L^\infty((-T,T))}<+\infty$ and as well, if $k\geq 1$, $||(v^\pm_\lambda)^{(k)}||_{L^2(\R)} = ||(v^\pm_\lambda)^{(k)}||_{L^2((-T,T))}<+\infty$.

   We shall approximate $u$, both in measure and in energy, by gluing together rescaled versions of the approximate optimal profiles $v_j\coloneqq v^{s_j}_\lambda$ (from now on we consider $\lambda$ fixed and drop the explicit dependence).  For this purpose, we partition $(0,1)$ into sub-intervals symmetric with respect to the jump-points, namely
    \begin{align*}
        I_j &\coloneqq  (b_{j-1}, b_j), \quad j=1,\ldots, M, \\
        b_j &\coloneqq \frac{t_j+t_{j+1}}{2}, \quad j=1,\ldots, M-1, \,b_0\coloneqq 0, \,
        b_{M}\coloneqq 1.
    \end{align*}
     Then, we define the sequences
    \begin{equation}
        \label{eq:recovery_sequence}
        u_\e(x) = \begin{cases}
            \displaystyle\sum_{j=1}^M v_{j}\Big( \frac{\lambda}{ \del}(x-t_j^\delta) \Big)\mathsf{1}_{I_j}(x), &\text{ for }\lambda\in (0,+\infty),\\
            \displaystyle\sum_{j=1}^M v_{j}\Big(\mkern1mu\frac{x-t_j^\delta}{ \e} \mkern1mu\Big)\mathsf{1}_{I_j}(x), &\text{ for }\lambda\in \{0,+\infty\},\\
        \end{cases}
    \end{equation}
    where $\mathsf{1}$'s denote indicator functions and the points $t^\delta_j$ satisfy $|t_j-t_j^\delta|\le\del$, but their precise value will be determined later with respect to $\lambda$.
    
    For $\e$ small enough, precisely such that 
    \begin{equation*}
        \tau_\e\coloneqq \e T+\delta<\tau,
    \end{equation*} $u_\e$ belongs to $H^{k+s}((0,1))$.
    Indeed, for sure $u_\e\in H^{k+s}(I_j)$ for each $j$. Furthermore, as $v_j\in H^{k+s}_{\rm loc}(\R)$, $u_\e\in H^k((0,1))$ by construction.  % and the construction is local, we have
    %\[\|u_\e\|^2_{H^k((0,1))} = \sum_{j=1}^M \|v_j\|^2_{H^k(I_j)}<+\infty.\]
    Then $u_\e\in H^{k+s}((0,1))$ as a consequence of the following claim:
    \begin{equation}
    \label{bound interactions}
        \int_{I_i}\int_{I_j} \frac{|u^{(k)}_\e(y)-u^{(k)}_\e(x)|^2}{|x-y|^{1+2s}} \,\de x\,\de y \leq \frac{C}{(\tau-\tau_\e)^{1+2s}}\big(\|u_\e^{(k)}\|^2_{L^2(I_j)} +  \|u_\e^{(k)}\|^2_{L^2(I_i)}\big).
    \end{equation}
    In order to prove the claim, note that the integrand can be non-zero only if
$|x-y|\geq \tau-\tau_\e$. Indeed, $|x-y|\geq 2\tau$  when  $|i-j|>1$, while, if $i=j+1$, 
\[u^{(k)}_\e(x)-u^{(k)}(y) \equiv 0,\]
for any $x\in\big(t_j^\delta+\e T, b_j)$ and $y\in(b_j,t_{j+1}^\delta-\e T\big)$. Thus, the integrand is non-zero if
\[|x-y| \geq \tau-\e T-\delta\geq \tau-\tau_\e,\]
and consequently 
\begin{align*}
   \int_{I_i}\int_{I_j} \frac{|u^{(k)}_\e(y)-u^{(k)}_\e(x)|^2}{|x-y|^{1+2s}} \,\de x\,\de y&\leq C\int_{I_i}\int_{I_j} \frac{|u^{(k)}_\e(y)-u^{(k)}_\e(x)|^2}{(\tau-\tau_\e)^{1+2s}} \,\de x\,\de y \\
    &\leq \frac{C}{(\tau-\tau_\e)^{1+2s}}\big(\|u_\e^{(k)}\|^2_{L^2(I_j)} +  \|u_\e^{(k)}\|^2_{L^2(I_i)}\big).
\end{align*}


    We note that, for $\e$ small enough, $u_\e$ may differ from $u$ only inside the intervals $t_j+(-\tau_\e,\tau_\e)$
    of total length $2\tau_\e\rightarrow 0$, thus proving convergence in measure. As for the energy, from \eqref{bound interactions}, since $\tau_\e\to 0$, we have:
\begin{equation}\label{eq:energy_difference_recovery}
    \begin{split}
          & F_{\e,\del}(u_\e) - \sum_{j=1}^M F_{\e,\del}(u_\e, I_j) \\
           &\qquad =\sum_{i\neq j} \e^{2(k+s)-1} \int_{I_i}\int_{I_j}a\Big(\frac{x}{\delta},\frac{y}{\delta}\Big)\frac{\big|{u^{(k)}_\e(y)-u^{(k)}_\e(x)}\big|^2}{|y-x|^{1+2s}} \,\de x\,\de y\\
           &\qquad\leq C \e^{2(k+s)-1}\big(\|u_\e^{(k)}\|^2_{L^2(I_j)} +  \|u_\e^{(k)}\|^2_{L^2(I_i)}\big)
    \end{split}
\end{equation}

If $k\geq 1$, we have by construction (see \eqref{eq:recovery_sequence}) that for any $j\in\{1,\dots, M\}$
\[\|u_\e^{(k)}\|^2_{L^2(I_j)} \leq C \e^{1-2k}\|v_j^{(k)}\|^2_{L^2(\R)},\] 
hence, the right-hand side of \eqref{eq:energy_difference_recovery} tends to 0 (as $O(\e^{2s})$). The only subtlety here is that $\e$ must be small enough if $\lambda \in (0,+\infty)$, say $\delta/\e \leq 2\lambda$.
    
If $k=0$, we note that $||u_\e||_{L^\infty(I_j)}\leq ||v_j||_{L^\infty(\R)}<+\infty$, that is, thus,
% we recall again Theorem \ref{coercivity}, which assures us that the $v_\e$, and hence the $u_\e$ are equibounded, thus
\[
\begin{split}
    \int_{I_j}|u_\e(x)|^2\,\de x \leq C |I_j|\leq C.
\end{split}
\]
and, consequently, the right-hand side of \eqref{eq:energy_difference_recovery} tends to zero (as $O(\e^{2s-1})$). These arguments prove that
    \[\limsup_{\substack{\e\to 0 \\ \del/\e\to\lambda}} F_{\e,\delta}(u_\e) \leq \sum_{j=1}^M
   \limsup_{\substack{\e\to 0 \\ \del/\e\to\lambda}}F_{\e,\delta}(u_\e,I_j),\]
   for every $\lambda$. To conclude, from (\ref{eq:approx_choice}) and the arbitrariness of $N$, it is sufficient to show that $$\limsup_{\e \rightarrow 0} F_{\e,\delta}(u_\e, I_j) \leq \Phi^{a_\lambda}(v_j).$$ 


    
    \paragraph{Case $\lambda\in (0,+\infty)$.}  We choose  $t_j^\delta = \delta \left\lfloor\frac{t_j}{\delta}\right\rfloor$, so that by the periodicity of $a$ we have $a\Big( \frac{x}{\lambda}+\frac{t^\delta_j}{\del}, \frac{y}{\lambda}+\frac{t^\delta_j}{\del}\Big) = a\left( \frac{x}{\lambda}, \frac{y}{\lambda}\right)$, thus, by a change of variable we get
    \begin{equation*}
      \begin{split}
             F_{\e,\delta}(u_\e, I_j)
             &\le\frac{\del}{\lambda\e}\int_\R W\big(v_j(  x)\big) \,\de  x\\
               &\qquad +\Big(\frac{\del}{\lambda\e}\Big)^{1-2(k+s)}\iint_{\R^2}a\Big( \frac{x}{\lambda}, \frac{y}{\lambda}\Big)\frac{|v_{j}^{(k)}(  y)-v^{(k)}_{j}( x)|^2}{|y-x|^{1+2s}} \,\de x\,\de y\\
                   &=\Phi^{a(\cdot/\lambda)}(v_j)+\mr o_\e(1).
      \end{split}
    \end{equation*}
    %where the last passage follows from the fact that $\delta/\e\to \lambda$.
    \paragraph{Case $\lambda = 0$.}     
    We choose $t_j^\delta$ as in the previous case, and we get
     \begin{equation*}
      \begin{split}
             F_{\e,\delta}&(u_\e, I_j)\\
             %&\leq \int_\R W\big(v_j(  x)\big) \,\de  x +\iint_{\R^2}a\Big( \frac{\e}{\del}x+\frac{t^\delta_j}{\del}, \frac{\e}{\del}y+\frac{t^\delta_j}{\del}\Big)\frac{|v_{j}^{(k)}(  y)-v^{(k)}_{j}( x)|^2}{|y-x|^{1+2s}} \,\de x\,\de y\\ 
              &\leq\int_\R W\big(v_j(  x)\big) \,\de  x +\iint_{\R^2}a\Big( \frac{\e}{\del}x, \frac{\e}{\del}y \Big)\frac{|v_{j}^{(k)}(  y)-v^{(k)}_{j}( x)|^2}{|y-x|^{1+2s}} \,\de x\,\de y\\ 
                   &=\Phi^{\bar a}(v_j)+\mr o_\e(1),
      \end{split}
    \end{equation*}
where the last passage follows from the Riemann--Lebesgue Lemma:
$$a\Big(\frac{\e}{\delta}\cdot\,,\frac{\e}{\delta}\cdot\Big)\overset{*}{\rightharpoonup}  \bar a=\int_0^1\int_0^1a(x,y)\,\de x\,\de y, \quad\text{in } L^\infty.$$
  
    
    \paragraph{Case $\lambda = +\infty$.}
    Here we assume $a$ to be continuous, so that  
    \[a_{\inf} = \min_{0\leq t\leq 1} a(t,t)=a(r,r)\] 
    is well defined and attained by some $r\in [0,1)$. We chose
    \[t^\delta_j = \del\Big(\Big\lfloor\frac{t_j}{\delta}-r\Big\rfloor + r\Big),\] so that, taking into account the periodicity of $a$,  
    \[a\Big( \frac{\e}{\del}x+\frac{t^\delta_j}{\del}, \frac{\e}{\del}y+\frac{t^\delta_j}{\del}\Big)= a\left(r+ \frac{\e}{\delta}x, r+ \frac{\e}{\delta}y\right).\]
    Since $\frac{\e}{\delta} \rightarrow 0$ and $a$ is continuous, it follows that pointwise
    \[a\left(r+ \frac{\e}{\delta}x, r+ \frac{\e}{\delta}y\right) \rightarrow a(r,r) = a_{\inf},\]
    and, by the Dominated Convergence Theorem,
    \[a \left(r+ \frac{\e}{\delta}\cdot\, , r+ \frac{\e}{\delta}\cdot\right) \overset{*}{\rightharpoonup}  a_{\inf}\]
    in the weak-$*$ convergence of $L^\infty$. Then, we can repeat the argument of the previous case.
\end{proof}

\end{document}
